% Work-MAT | 2. rocnik | pripravovana hodina 74
% Návaznost: uživatelské připravy.zip (hodiny 01–20), autorská učebnice na GitHubu.
% Realisticky.cz: didaktická inspirace v posloupnosti a typech úloh; zadání jsou vlastní.
\documentclass[11pt,a4paper]{article}
\input{../../sablony/pripravy/priprava.sty}
\nastavPripravu{2. ročník}{Trigonometrie}{74}{Opakování trigonometrie}
\begin{document}
\blok{Motivace a návaznost}
Krátké opakování před T6: kombinace sinové a kosinové věty, obvodů a obsahu. Žáci navrhnou postup a až poté dosazují hodnoty.
\poznamka{Orientační nabídka pro 45--55 minut: krátké zopakování a motivace, nový výklad, tři společné vzorové úlohy, samostatná práce. Vyber podle tempa třídy; zbývající zadání slouží jako rezerva.}
\blok{Výklad}
Vedle výsledků hodnotíme úplnost náčrtku, správné označení stran a úhlů a přiměřené zaokrouhlení.
\blok{Společné příklady}
\begin{prikladReseniBox}[title={Příklad 1}]
Vyřeš, resp. urči: $a=3,b=4,\gamma=90^\circ:\;c,S$.
\tcblower
\[\begin{aligned}c=\sqrt{9+16}=5 \\ S=\frac{1}{2}\cdot3\cdot4=6\end{aligned}\]
\end{prikladReseniBox}
\begin{prikladReseniBox}[title={Příklad 2}]
Vyřeš, resp. urči: $a=6,\alpha=30^\circ,\beta=45^\circ:\;b$.
\tcblower
\[\begin{aligned}b=6\sqrt2\end{aligned}\]
\end{prikladReseniBox}
\poznamka{Nech jednotlivé žáky obhajovat zvolený vzorec; výsledky kontrolujte výpočtem nebo odhadem.}
\blok{Samostatné procvičování}
\begin{zadaniReseniBox}[title={Samostatná práce 1}]
\begin{enumerate}[label=\alph*)]
\item $a=8,b=5,\gamma=30^\circ:\;S$.
\item $a=6,b=4,\gamma=90^\circ:\;S$.
\item $a=10,\alpha=30^\circ,\beta=60^\circ:\;b$.
\item $a=4,\alpha=30^\circ,\beta=90^\circ:\;b$.
\item $a=5,b=5,\gamma=60^\circ:\;c$.
\item $a=6,b=8,\gamma=90^\circ:\;c$.
\end{enumerate}
\tcblower
\begin{enumerate}[label=\alph*)]
\item $10$.
\item $12$.
\item $10\sqrt3$.
\item $8$.
\item $5$.
\item $10$.
\end{enumerate}
\end{zadaniReseniBox}
\blok{Dobrovolné rozšíření}
\begin{zadaniReseniBox}[title={Rozšíření 1}]
\begin{enumerate}[label=\alph*)]
\item Nakresli vlastní realistickou technickou úlohu, označ známé údaje a navrhni řešení..
\end{enumerate}
\tcblower
\begin{enumerate}[label=\alph*)]
\item Ověř náčrtek, označení a smysluplnost jednotek..
\end{enumerate}
\end{zadaniReseniBox}
\blok{Úkol na doma}
U každé úlohy načrtni a označ trojúhelník, napiš použitou větu a příslušné jednotky. Tam, kde je to účelné, proveď odhad výsledku.
\begin{zadaniReseniBox}[title={Domácí úkol}]
\begin{enumerate}[label=\alph*)]
\item $a=4,b=6,\gamma=60^\circ:\;S$.
\item $a=8,b=3,\gamma=30^\circ:\;S$.
\item $a=7,b=7,\gamma=120^\circ:\;c$.
\item $a=5,b=12,\gamma=90^\circ:\;c$.
\item $a=5,\alpha=30^\circ,\beta=30^\circ:\;b$.
\item $a=10,b=10,\gamma=60^\circ:\;S$.
\item $a=6,\alpha=30^\circ,\beta=45^\circ:\;b$.
\item $a=8,\alpha=30^\circ,\beta=90^\circ:\;b$.
\end{enumerate}
\tcblower
\begin{enumerate}[label=\alph*)]
\item $6\sqrt3$.
\item $6$.
\item $7\sqrt3$.
\item $13$.
\item $5$.
\item $25\sqrt3$.
\item $6\sqrt2$.
\item $16$.
\end{enumerate}
\end{zadaniReseniBox}
\poznamka{Domácí úkol obsahuje 8 položek; předpokládaná domácí práce přibližně 15--30 minut. Po větších blocích jsou vědomě promíchané dříve probrané typy.}
\end{document}
